\section{Reflective Statement}
\label{sec:reflection}
\subsection{MVP - Most Vital Papers}
The two publications that were critical to this review were Venter's thesis \cite{venterSOFTWAREDEFINEDPULSEDOPPLERRADAR}, and Hoefler and Ballis rules for scientific benchmarking \cite{hoeflerScientificBenchmarkingParallel2015a}.

Early in the literature reading, my critiques were largely instinctive. Paper after paper
that I read reported large speed-ups against a serial CPU, and my reaction to this was that a GPU beating a single-core CPU is the least I would expect! I could very easily point out what was wrong in these experiments, but struggled to formulate what good would look like for this. Venter gave me the reference point that I was searching for, he computed the arithmetic intensity of every kernel in a complete pulse-doppler chain, measured each against peak and copy-kernel bandwidth and timed transfers separately from computation. This work became the model for Table~\ref{tab:workload} and it refined my first gap from ``no one has characterised a radar chain'' to ``it has been done once, on a 2008 GPU with one receive channel''. Additionally, it taught me that an exemplar paper is not necessarily a standard, Venters thesis fails to report any CPU baselines and does not include spread in his measurements, which fails two of my own six criteria.

After reading through Venter's thesis, I began the literature search for answering the question of what good looked like in these measurements. Hoefler and Balli gave me the criteria that could be applied to any paper. Their rules on fairness confirmed what I was already suspecting when reading through the initial radar papers. Their analytical and statistical rules were new to me, and became the analysis-depth and timing-statistic criteria in Table~\ref{tab:rubric}. Their advice to judge latency-sensitive based on upper percentiles is something that I was already familiar with from the hard-deadlines imposed on radar processing, where a single late CPI has a cascading bottlenecking property.

Due to a lack of detail in the published papers, I could not apply all twelve rules highlighted in the benchmarking criteria. Their rules for comparing results required a statistical distribution, or at least a spread of results, which was not provided by the papers surveyed.

Applying both sources also exposed a tension in my own method. Venter, Garner, and Zaknoun contain the most detailed reviews, and all are theses. Although my own method prefers peer-reviewed work. The reason behind this, I think, is due to format: a thesis is assessed on its full method and has room to report it, whereas conference / journal entries are often subject to tight page limits so leave out the exact detail that Hoefler and Balli advocate for. In the process of cutting this review to the five-page limit, I found the same trade-off of removing detail to fit the page count.


\subsection{Use of Generative AI}

I made use of Anthropics Claude throughout the lifecycle of this review, so this reflection will be stage-by-stage.

I first used it as a research assistant to build and screen the corpus of papers. Rather than letting it work from general knowledge and its own context window, I connected to my Obsidian vault which serves as my second brain for my PhD, this allowed all of its summaries / insights to be derived from ground truths, in hope of minimising the risk of hallucinations and errors. While it did reduce the amount of error, it did not eliminate it. My first corpus came from AI web-searching and manual snowballing, but when I came to write Section~\ref{sec:method}, I found that it had not recorded hit counts, so it was not reproducable. To fix this, I had to perform a manual database search and, for truth-verification, performed a manual screening of all the papers it had returned. Doing so, I discoverd an additional 3 relevant papers to the review that were missed by AI.

Working as a part-time EngD student with about two days a week for study, I used Claude (integrated with my Notion project-management tool) to help decompose the coursework specification into a series of clear tasks and roadmapped the completion, giving myself deadlines each week to ensure I meet the deadline alongside other course commitments.

Once my initial searches were complete, Claude extracted details from each paper and generated them into an evidence matrix to compare the studies. I manually verified every row against the full text, and the errors I found were invented facts that the papers did not mention. These were fairly minor reasoning errors, such as counting in the wrong units.

Finally, Claude helped produce the skeleton of this report writing out a draft that I could expand upon. This is where connecting Claude to my Obsidian vault proved most useful, as it could produce the plan based on the evidence vault for my PhD. Additionally, in the writing stage, it helped with the specific wording of some claims. Despite helping in the planning and early drafting, I still decided what would be included in the report and what would be cut in favour of the page count. I found minor 'hallucinations' in the produced plan which needed corrected.

In conclusion, I would use generative AI again for aiding in the planning, performing a first pass over a large literature base, but not for deciding what the evidence supports or for drawing any conclusions between papers. Personally, I would not trust any figure that was derived from an LLM until I had recomputed it myself. While AI made the process of doing this review faster, and allowed me to focus on the most-important parts, while semi-automating some of the management, every claim in the review still rests in my own reading and judgment.